\documentclass[11pt,letterpaper,landscape]{article}

\usepackage[
  top=0.55cm,
  bottom=0.55cm,
  left=0.65cm,
  right=0.65cm
]{geometry}
\usepackage{xcolor}
\usepackage{fontspec}
\usepackage{tcolorbox}
\tcbuselibrary{skins}
\usepackage{tabularx}
\usepackage{listings}
\usepackage{amsmath}

\setmainfont{Libertinus Sans}
\setmonofont{Libertinus Mono}
\pagestyle{empty}
\setlength{\parindent}{0pt}
\setlength{\tabcolsep}{4pt}
\renewcommand{\arraystretch}{1.12}
\newcolumntype{Y}{>{\raggedright\arraybackslash}X}

\definecolor{tealbox}{HTML}{079A9A}
\definecolor{orangeTitle}{HTML}{FF5A00}

\lstdefinestyle{pyinline}{
  basicstyle=\ttfamily\footnotesize,
  columns=fullflexible,
  keepspaces=true,
  breaklines=true,
  showstringspaces=false
}
\lstset{style=pyinline}

\newcommand{\code}[1]{\texttt{\detokenize{#1}}}
\newcommand{\entry}[2]{\textbf{\code{#1}} & #2\\}
\newcommand{\smallentry}[2]{\code{#1} & #2\\}
\newcommand{\colw}{0.318\textwidth}

\newtcolorbox{cheatbox}[1]{
  enhanced,
  sharp corners,
  boxrule=0pt,
  colback=white,
  left=2.5mm,
  right=2.5mm,
  top=1.2mm,
  bottom=1.2mm,
  before skip=2.5mm,
  after skip=2.5mm,
  title=#1,
  coltitle=white,
  fonttitle=\bfseries\large,
  colbacktitle=tealbox,
  attach boxed title to top center={yshift=-1mm},
  boxed title style={sharp corners, boxrule=0pt, colback=tealbox, left=3mm, right=3mm},
}

\newcommand{\cheattitle}{
  \begin{center}
    {\fontsize{23}{25}\selectfont\bfseries\color{orangeTitle} Python para Solemne 3 - Cheat Sheet}
  \end{center}
  \vspace{-3mm}
}

\begin{document}
\cheattitle

\noindent
\begin{minipage}[t]{\colw}
\vspace{0pt}
\begin{cheatbox}{Pandas: crear e inspeccionar}
\begin{lstlisting}
import numpy as np
import pandas as pd

df = pd.DataFrame(datos)
print(df.head())
print(df.shape)
print(df.dtypes)
print(df.isna().sum())
\end{lstlisting}
\begin{tabularx}{\linewidth}{@{}lY@{}}
\entry{df.head()}{primeras filas}
\entry{df.shape}{filas y columnas}
\entry{df.dtypes}{tipo de cada columna}
\entry{df.isna().sum()}{nulos por columna}
\end{tabularx}
\end{cheatbox}

\begin{cheatbox}{Pandas: limpiar datos}
\begin{lstlisting}
cols = ["temperatura_K", "radio_Rsol",
        "distancia_pc", "magnitud_V"]

df_limpio = df.dropna(subset=cols).copy()
\end{lstlisting}
\begin{tabularx}{\linewidth}{@{}lY@{}}
\entry{dropna(subset=cols)}{elimina filas con nulos en columnas clave}
\entry{copy()}{crea una copia para modificar sin ambiguedad}
\entry{len(df)}{cantidad de filas}
\end{tabularx}
\end{cheatbox}

\begin{cheatbox}{Pandas: columnas derivadas}
\begin{lstlisting}
df_limpio["luminosidad_rel"] = (
    df_limpio["radio_Rsol"]**2 *
    (df_limpio["temperatura_K"] / 5778)**4
)
\end{lstlisting}
Las operaciones con columnas se aplican fila a fila.
\end{cheatbox}
\end{minipage}
\hfill
\begin{minipage}[t]{\colw}
\vspace{0pt}
\begin{cheatbox}{Pandas: filtros}
\begin{lstlisting}
mask = (
    (df["distancia_pc"] < 20) &
    (df["magnitud_V"] < 0.5) &
    (df["luminosidad_rel"] > 5)
)

sub = df.loc[mask, ["nombre", "tipo"]]
\end{lstlisting}
\begin{tabularx}{\linewidth}{@{}lY@{}}
\entry{\&}{combinar condiciones}
\entry{|}{una condicion u otra}
\entry{df.loc[mask, cols]}{filas filtradas y columnas elegidas}
\entry{(...) \& (...)}{use parentesis en cada condicion}
\end{tabularx}
\end{cheatbox}

\begin{cheatbox}{Pandas: agrupar}
\begin{lstlisting}
resumen = df.groupby("tipo").agg({
    "nombre": "count",
    "temperatura_K": "mean",
    "distancia_pc": "median",
    "luminosidad_rel": "mean"
})

resumen = resumen.rename(columns={"nombre": "cantidad"})
resumen = resumen.sort_values("luminosidad_rel",
                              ascending=False)
\end{lstlisting}
\end{cheatbox}

\begin{cheatbox}{Pandas: graficar por grupo}
\begin{lstlisting}
for tipo in df["tipo"].unique():
    sub = df[df["tipo"] == tipo]
    plt.scatter(sub["temperatura_K"],
                sub["luminosidad_rel"],
                label=tipo)
plt.legend()
\end{lstlisting}
\end{cheatbox}
\end{minipage}
\hfill
\begin{minipage}[t]{\colw}
\vspace{0pt}
\begin{cheatbox}{Algoritmos: usar copia}
\begin{lstlisting}
periodos = [12.4, 3.1, 8.7]
ordenados = bubble_sort(periodos.copy())

print(periodos)    # original
print(ordenados)   # ordenada
\end{lstlisting}
\begin{tabularx}{\linewidth}{@{}lY@{}}
\entry{lista.copy()}{evita modificar la lista original}
\entry{ordenados[0]}{minimo si esta ordenada}
\entry{ordenados[-1]}{maximo si esta ordenada}
\entry{len(lista)}{cantidad de elementos}
\end{tabularx}
\end{cheatbox}

\begin{cheatbox}{Mediana con lista ordenada}
\begin{lstlisting}
n = len(ordenados)

if n % 2 == 1:
    mediana = ordenados[n//2]
else:
    mediana = (ordenados[n//2 - 1] +
               ordenados[n//2]) / 2
\end{lstlisting}
\end{cheatbox}

\begin{cheatbox}{Buscar e interpretar}
\begin{lstlisting}
pos = binary_search(ordenados, 14.2)

if pos == -1:
    print("no encontrado")
else:
    print("posicion:", pos)
\end{lstlisting}
\end{cheatbox}
\end{minipage}

\newpage
\cheattitle

\noindent
\begin{minipage}[t]{\colw}
\vspace{0pt}
\begin{cheatbox}{Bubble Sort}
\begin{lstlisting}
def bubble_sort(lista):
    n = len(lista)
    for i in range(n - 1):
        for j in range(n - 1 - i):
            if lista[j] > lista[j + 1]:
                lista[j], lista[j + 1] = (
                    lista[j + 1], lista[j]
                )
    return lista
\end{lstlisting}
Ordena comparando vecinos. Use una copia si quiere conservar la lista original.
\end{cheatbox}

\begin{cheatbox}{Binary Search}
\begin{lstlisting}
def binary_search(lista_ordenada, objetivo):
    low, high = 0, len(lista_ordenada) - 1
    while low <= high:
        mid = (low + high) // 2
        if lista_ordenada[mid] == objetivo:
            return mid
        elif objetivo < lista_ordenada[mid]:
            high = mid - 1
        else:
            low = mid + 1
    return -1
\end{lstlisting}
Solo es valida sobre una lista ordenada.
\end{cheatbox}
\end{minipage}
\hfill
\begin{minipage}[t]{\colw}
\vspace{0pt}
\begin{cheatbox}{Ajuste con \texttt{np.polyfit}}
\begin{lstlisting}
import numpy as np

a, b, c = np.polyfit(d, R, deg=2)
R_fit = a*d**2 + b*d + c
\end{lstlisting}
\begin{tabularx}{\linewidth}{@{}lY@{}}
\entry{deg=1}{modelo lineal: \code{m*x + b}}
\entry{deg=2}{modelo cuadratico: \code{a*x**2 + b*x + c}}
\entry{polyfit}{retorna coeficientes desde mayor grado}
\end{tabularx}
\end{cheatbox}

\begin{cheatbox}{Residuos y \(R^2\)}
\begin{lstlisting}
residuos = R - R_fit
SS_res = np.sum(residuos**2)
SS_tot = np.sum((R - np.mean(R))**2)
R2 = 1 - SS_res/SS_tot
\end{lstlisting}
\begin{tabularx}{\linewidth}{@{}lY@{}}
\entry{residuos}{dato medido menos modelo}
\entry{R2 cercano a 1}{ajuste explica bien los datos}
\end{tabularx}
\end{cheatbox}

\begin{cheatbox}{Graficar ajuste}
\begin{lstlisting}
x_fino = np.linspace(0, 6, 200)
y_fino = a*x_fino**2 + b*x_fino + c

plt.scatter(d, R, label="datos")
plt.plot(x_fino, y_fino, label="ajuste")
plt.legend()
plt.grid(True)
\end{lstlisting}
\end{cheatbox}
\end{minipage}
\hfill
\begin{minipage}[t]{\colw}
\vspace{0pt}
\begin{cheatbox}{Raices con \texttt{np.roots}}
\begin{lstlisting}
# Resolver a*x**2 + b*x + c = 0
raices = np.roots([a, b, c])

# Resolver a*d**2 + b*d + c = 12
raices = np.roots([a, b, c - 12])
\end{lstlisting}
\begin{tabularx}{\linewidth}{@{}lY@{}}
\entry{[a,b,c]}{coeficientes en orden descendente}
\entry{c - 12}{mover el umbral al lado izquierdo}
\entry{np.roots}{puede entregar raices complejas}
\end{tabularx}
\end{cheatbox}

\begin{cheatbox}{Filtrar raices fisicas}
\begin{lstlisting}
soluciones = []
for r in raices:
    if abs(r.imag) < 1e-10 and 0 <= r.real <= 6:
        soluciones.append(r.real)
\end{lstlisting}
\begin{tabularx}{\linewidth}{@{}lY@{}}
\entry{r.imag}{parte imaginaria}
\entry{r.real}{parte real}
\entry{0 <= r.real <= 6}{rango fisico o medido}
\end{tabularx}
\end{cheatbox}

\begin{cheatbox}{Errores comunes}
\begin{tabularx}{\linewidth}{@{}lY@{}}
\smallentry{and}{en pandas use \code{\&} con parentesis}
\smallentry{binary_search(lista)}{debe ser lista ordenada}
\smallentry{np.roots([c,b,a])}{orden incorrecto}
\smallentry{tomar toda raiz}{filtre complejas y no fisicas}
\end{tabularx}
\end{cheatbox}
\end{minipage}

\end{document}
